% Persian -- what the reading editions' preamble leaves to the language.
% Input by lib/frank-preamble.tex as \input{\FrankLib/lang/\BookLang.tex};
% docs/languages.md section 6 lists the hooks every file in this directory
% must define, and this one is the model the others were copied from.
%
% This file is where the Persian-only preamble went when languages arrived.
% Nothing in it was new: the babel locale, the fonts, the harakat filter, the
% digit map and the how-to-read text are the ones the first editions were
% built with, and the Persian book rebuilt to the same PDF text (the
% pdftotext before/after test that build.sh's history relies on).  The one
% thing added since is two sentences on the how-to page, saying what the
% three forms of a verb cannot: that داشتن's present takes no mi-, and why
% a verb with a preverb is hyphenated.  The chapters print as they did.

% ---- the babel language and the switches ----------------------------------
\newcommand{\FrankLangName}{persian}
\FrankRTLtrue            % chunk right, gloss left; \pw isolates with \babelsublr
\FrankHasBaretrue        % pass 3 is the same text without the harakat
\FrankHasAlttrue         % pass 4 is a FONT pass: the same text in nastaliq
\FrankHasReadingfalse    % no kana-like reading line
\FrankHasVertfalse

% ---- fonts ------------------------------------------------------------------
% The faces come from the registry's "tex" record (lib/languages.json), so no
% font is named here: \FrankPickFont walks tex.main then tex.main_fallbacks and
% takes the first face installed.  Vazirmatn and the nastaliq travel with the
% toolbox (lib/fonts, on OSFONTDIR), so on this machine the pick is the same
% as the names the preamble once carried literally.
\babelprovide[import=fa,onchar=ids fonts]{persian}
\FrankPickFont{\FrankMainFontName}{fa}{main}{main_fallbacks}
\babelfont[persian]{rm}[Script=Arabic,Language=Parsi,Renderer=HarfBuzz]{\FrankMainFontName}
% The alternate face: nastaliq, the hand Persian poetry is set in.  \nastaliqfont
% is the historical name and stays usable; \FrankAltFont is what the core uses.
\FrankPickFont{\FrankAltFontName}{fa}{alt}{alt_fallbacks}
\newfontfamily\nastaliqfont{\FrankAltFontName}%
  [Script=Arabic,Language=Urdu,Renderer=HarfBuzz]
\newcommand{\FrankAltFont}{\nastaliqfont}

% ---- the vocabulary line -----------------------------------------------------
% \vb{infinitive}{sound}{present stem}{sound}{past stem}{sound}{meaning}: these
% two are the labels printed before the second and third forms.  They are the
% core preamble's own defaults -- Persian is what it was written for -- and are
% spelt out here because the reader takes the same pair from the registry
% (lib/languages.json, "vb_labels": ["pres.", "past"], beside "vb_forms",
% which names the three forms in words for the chunk editor); the two must
% agree, or the PDF and the reader of one book label the same stem
% differently.  lib/newlang.py --check compares them.
%
% Every pair is filled in Persian: every verb has both stems, and the present
% one is the one nobody can guess (دیدن, بین).  What the three forms cannot
% say goes in ONE parenthesis after the meaning, and for Persian that is a
% closed list of one: داشتن, whose present takes no mi- -- دارم, never
% می‌دارم -- is \vb{داشتن}{dāštan}{دار}{dār}{داشت}{dāšt}{to have (pres.
% without mi-)}.  بودن keeps باش as its stem and هست stays a \dw of its own
% (check_batch refuses \vb{بودن}{...}{هست}).  A verb with a preverb is
% hyphenated in all three sounds, bar-gaštan, bar-gard, bar-gašt, because
% mi-, be- and na- go in after the preverb: bar-mi-gardam.  A compound is the
% light verb's \vb with an empty meaning and a \bw for the word it carries.
% The spoken present of a video (coll. می‌گم mi-gam for گفتن) is the video's
% plain line only; a book gives the written stems and nothing else.
\newcommand{\FrankVbPres}{pres.}
\newcommand{\FrankVbPast}{past}

% ---- the two Lua filters ---------------------------------------------------
% frank_strip: pass 1 carries the short vowels, the way a children's edition
% does.  Passes 3 and 4 must not, but the Persian is only written once -- so
% the harakat (U+064B..U+0652, the registry's `strip`) are filtered out of the
% copy that feeds them.
% frank_latin: Persian digits are the label, Latin digits are the id, so a
% destination name has to be transliterated (U+06F0..U+06F9 -> ASCII).
% Both are called through \directlua, which is expandable and so survives an
% \edef.  Keep the bodies free of #, % and a literal backslash (NOTES section 9).
\directlua{
  function frank_strip(s)
    local t, n = {}, 0
    for _, c in utf8.codes(s) do
      if c < 0x064B or c > 0x0652 then n = n + 1 ; t[n] = utf8.char(c) end
    end
    return table.concat(t)
  end
  function frank_latin(s)
    local t, n = {}, 0
    for _, c in utf8.codes(s) do
      if c >= 0x06F0 and c <= 0x06F9 then c = c - 0x06F0 + 0x30 end
      n = n + 1 ; t[n] = utf8.char(c)
    end
    return table.concat(t)
  end
}

% ---- the "How to read this" page -----------------------------------------------
% Printed by lib/frank-frontmatter.tex under its heading.  \pw is defined by the
% core preamble, after this file is read; it is only expanded when the page is set.
\newcommand{\FrankHowTo}{%
Every passage comes four times.

\smallskip
\textbf{First} the whole sentence in naskh, with the short vowels written in
as in a children's edition, and nothing else. Try to read it. Get as far as
you can and do not worry about what you cannot get — the point is the attempt,
made before any help arrives.

\smallskip
\textbf{Second} the same sentence in small chunks, Persian on the right, and
on the left three lines: how it sounds, then how to look each word up, then
what it means. Now you find out how much you had right.

\smallskip
\textbf{Third} the same passage again in naskh, this time without the vowels,
the way Persian is really written. You have read it twice, so it should carry
you.

\smallskip
\textbf{Fourth} the same again in nastaliq, the hand Persian poetry is set in.
You already know what it says, which is the only comfortable way to learn a
new hand.

\smallskip
Verbs are given as infinitive, present stem, past stem, meaning — the present
stem is the one you cannot guess. Where a verb is a compound made with
\pw{کردن}, \pw{شدن}, \pw{داشتن} or \pw{بردن}, only the meaning of the word
carried by it is given. What the three forms cannot tell you stands in
brackets after the meaning: the present of \pw{داشتن} takes no \textit{mi-}
(\pw{دارم} \textit{d\=aram}), and its entry says so. A verb that begins with
a prefix is hyphenated after it, \textit{bar-gard}, because that is where
\textit{mi-} goes in: \textit{bar-mi-gardam}.

\smallskip
\textit{-e} after a noun is the \textit{ez\=afe}, the vowel that links it to
what follows; it is written here as a kasra but is invisible in ordinary
Persian. \textit{\=a} is the long open vowel of \textit{f\=ar}, \textit{x} the
\textit{ch} of \textit{loch}, \textit{\v{s}} is \textit{sh}, \textit{\v{c}} is
\textit{ch}, \textit{\v{z}} the \textit{s} of \textit{measure}, \textit{'} a
glottal catch.
}
